\documentclass[aspectratio=169]{beamer}

% ==== Theme & basic look ====
\usetheme{Madrid}
\usecolortheme{seahorse}
\setbeamertemplate{navigation symbols}{}
\setbeamertemplate{footline}[frame number]


% Packages for mathematics  
\usepackage{amsmath}  
\usepackage{amsfonts}  
\usepackage{amssymb}  
\usepackage{CJKutf8}
% Support for \Tilde macro used throughout
\providecommand{\Tilde}[1]{\tilde{#1}}
\providecommand{\proofname}{Proof}
\newenvironment<>{proofs}[1][\proofname]{%
    \par
    \def\insertproofname{#1.}%
    \usebeamertemplate{proof begin}#2}
  {\usebeamertemplate{proof end}}
\makeatother

% Math shorthand and vector symbols
\newcommand{\grad}{\nabla}
\newcommand{\EE}{\mathbb{E}}
\newcommand{\dd}{\mathrm{d}}
\newcommand{\kbT}{k_B T}
\newcommand{\vx}{\mathbf{x}}
\newcommand{\va}{\mathbf{a}}
\newcommand{\vs}{\mathbf{s}}
\newcommand{\vr}{\mathbf{r}}
\newcommand{\vz}{\mathbf{z}}

% Bold vector/time-indexed variables
\newcommand{\vst}{\mathbf{s}_t}
\newcommand{\vat}{\mathbf{a}_t}
\newcommand{\vrt}{\mathbf{r}_t}
\newcommand{\vzt}{\mathbf{z}_t}
\newcommand{\vstp}{\mathbf{s}_{t+1}}
\newcommand{\vrtp}{\mathbf{r}_{t+1}}

% Title page details:   
\title[Understanding World Models]{Understanding World Models in Machine Learning and AGI}  
\author{Pipi Hu}  
\institute[BIMSA]{Beijing Institute of Mathematical Sciences and Applications}  
\date{\today}  

\pgfdeclareimage[width=2cm]{logo}{figs/logo.png}
\logo{\pgfuseimage{logo}}

% Outline slide (moved to later; avoid frame before \begin{document})

\setbeamertemplate{navigation symbols}{}
\AtBeginSection[]  
{  
    \begin{frame}<beamer>{Outline}         
        \tableofcontents[currentsection]  
    \end{frame}  
}  
  
\AtBeginSubsection[]  
{  
    \begin{frame}<beamer>{Outline}         
        \tableofcontents[currentsection,currentsubsection]  
    \end{frame}  
}  

\begin{document}  
  
% Title slide  
\begin{frame}  
  \titlepage  
\end{frame}  

\begin{frame}{Feynman's words}
\centering
    What I cannot create, I do not understand.
\end{frame}
  
% Presentation slides  

\section{Introduction: What are World Models?}

\begin{frame}{The Concept of World Models}
A \textbf{world model} is an internal representation that an AI system constructs to simulate and understand its environment.

\begin{itemize}
    \item \textcolor{red}{Abstraction}: Extract high-level features from raw sensory data
    \item \textcolor{red}{Prediction}: Forecast future states and rewards
    \item \textcolor{red}{Simulation}: Internally simulate virtual scenarios
    \item \textcolor{red}{Interaction}: Infer intentions and beliefs of other agents
\end{itemize}

\textbf{Key Insight}: World models enable AI systems to reason about the world without direct interaction, crucial for achieving AGI.
\end{frame}

\begin{frame}{Historical Context: From Reinforcement Learning to AGI}
\begin{enumerate}
    \item \textbf{Classical RL}: Direct interaction with environment
    \begin{equation}
        \pi^*(s) = \arg\max_a \EE\left[\sum_{t=0}^{\infty} \gamma^t r_{t+1} \mid s_0 = s, a_0 = a\right]
    \end{equation}
    \item \textbf{Model-based RL}: Learn environment dynamics
    \begin{equation}
        \hat{p}(s_{t+1}, r_{t+1} \mid s_t, a_t) \approx p(s_{t+1}, r_{t+1} \mid s_t, a_t)
    \end{equation}
    \item \textbf{World Models}: Internal simulation and planning
    \begin{equation}
        \text{World Model} = \{\text{Encoder}, \text{Dynamics}, \text{Decoder}\}
    \end{equation}
\end{enumerate}
\end{frame}

\section{Mathematical Foundations of World Models}

\begin{frame}{Formal Definition of World Models}
A world model $\mathcal{W}$ consists of three components:

\begin{enumerate}
    \item \textbf{Encoder}: $z_t = \text{Enc}(s_t)$ - maps observations to latent states
    \item \textbf{Dynamics}: $z_{t+1} = f(z_t, a_t)$ - predicts next latent state
    \item \textbf{Decoder}: $\hat{s}_t = \text{Dec}(z_t)$ - reconstructs observations
\end{enumerate}

\textbf{Objective}: Learn a compressed representation that captures the essential dynamics of the environment.
\end{frame}

\begin{frame}{The World Model Learning Problem}
Given a sequence of observations and actions $\{(s_t, a_t)\}_{t=1}^T$, we want to learn:

\begin{equation}
\min_{\theta} \mathcal{L}_{\text{world}} = \EE\left[\sum_{t=1}^T \|\hat{s}_t - s_t\|^2 + \lambda \mathcal{L}_{\text{consistency}}\right]
\end{equation}

where $\mathcal{L}_{\text{consistency}}$ ensures the learned dynamics are consistent with the true environment dynamics.

\textbf{Key Challenge}: The dynamics model must be both accurate and generalizable to unseen scenarios.
\end{frame}

\begin{frame}{Variational World Models}
Following the variational autoencoder framework, we can formulate world models as:

\begin{equation}
p(z_{1:T} \mid s_{1:T}, a_{1:T}) = \prod_{t=1}^T p(z_t \mid z_{t-1}, a_{t-1})
\end{equation}

\textbf{Variational Lower Bound}:
\begin{equation}
\log p(s_{1:T} \mid a_{1:T}) \geq \EE_{q(z_{1:T})}\left[\sum_{t=1}^T \log p(s_t \mid z_t) - \text{KL}(q(z_t \mid z_{t-1}, a_{t-1}) \| p(z_t \mid z_{t-1}, a_{t-1}))\right]
\end{equation}

This formulation allows for uncertainty quantification and better generalization.
\end{frame}

\section{World Models in Reinforcement Learning}

\begin{frame}{Model-Based Reinforcement Learning with World Models}
Traditional model-based RL learns environment dynamics:
\begin{equation}
\hat{p}(s_{t+1}, r_{t+1} \mid s_t, a_t) \approx p(s_{t+1}, r_{t+1} \mid s_t, a_t)
\end{equation}

\textbf{World Model Approach}: Learn dynamics in latent space
\begin{equation}
\hat{z}_{t+1} = f_\theta(z_t, a_t), \quad \hat{r}_{t+1} = g_\theta(z_t, a_t)
\end{equation}

\textbf{Advantages}:
\begin{itemize}
    \item More efficient learning in high-dimensional spaces
    \item Better generalization to unseen states
    \item Enables planning in latent space
\end{itemize}
\end{frame}

\begin{frame}{Planning with World Models}
Once we have a world model, we can plan by:

\begin{enumerate}
    \item \textbf{Rollout Planning}: Simulate trajectories in latent space
    \begin{equation}
        z_{t+1} = f(z_t, a_t), \quad \hat{r}_{t+1} = g(z_t, a_t)
    \end{equation}
    \item \textbf{Value Function Learning}: Learn $V(z)$ in latent space
    \begin{equation}
        V(z_t) = \EE\left[\sum_{k=0}^H \gamma^k r_{t+k} \mid z_t\right]
    \end{equation}
    \item \textbf{Policy Optimization}: Update policy based on simulated experience
\end{enumerate}

\textbf{Key Insight}: Planning in compressed latent space is computationally more efficient than planning in raw observation space.
\end{frame}

\begin{frame}{Dreamer Algorithm: A Case Study}
The Dreamer algorithm demonstrates world model learning:

\begin{enumerate}
    \item \textbf{Representation Learning}: 
    \begin{equation}
        z_t \sim q_\phi(z_t \mid s_t, z_{t-1}, a_{t-1})
    \end{equation}
    \item \textbf{Dynamics Learning}:
    \begin{equation}
        \hat{z}_{t+1} \sim p_\theta(z_{t+1} \mid z_t, a_t)
    \end{equation}
    \item \textbf{Reward Prediction}:
    \begin{equation}
        \hat{r}_{t+1} = r_\theta(z_t, a_t)
    \end{equation}
    \item \textbf{Policy Learning}: Learn policy in latent space
    \begin{equation}
        \pi_\psi(a_t \mid z_t) = \text{softmax}(Q_\psi(z_t, a_t))
    \end{equation}
\end{enumerate}
\end{frame}

\section{World Models for AGI}

\begin{frame}{The Role of World Models in AGI}
World models are fundamental for achieving AGI because they provide:

\begin{enumerate}
    \item \textbf{Generalization}: Transfer knowledge across domains
    \item \textbf{Abstraction}: Extract high-level concepts
    \item \textbf{Planning}: Reason about future scenarios
    \item \textbf{Imagination}: Generate novel scenarios
\end{enumerate}

\textbf{AGI Requirements}:
\begin{equation}
\text{AGI} = \{\text{World Model}, \text{Reasoning}, \text{Learning}, \text{Planning}\}
\end{equation}
\end{frame}

\begin{frame}{Hierarchical World Models}
For complex environments, hierarchical world models are necessary:

\begin{equation}
\mathcal{W}_{\text{hier}} = \{\mathcal{W}_1, \mathcal{W}_2, \ldots, \mathcal{W}_L\}
\end{equation}

where each level $\mathcal{W}_i$ operates at different temporal and spatial scales.

\textbf{Multi-scale Dynamics}:
\begin{equation}
z_t^{(i)} = f^{(i)}(z_{t-1}^{(i)}, a_{t-1}^{(i)}, z_t^{(i-1)})
\end{equation}

This enables reasoning at multiple levels of abstraction, crucial for AGI.
\end{frame}

\begin{frame}{World Models and Theory of Mind}
For social AGI, world models must include models of other agents:

\begin{equation}
\mathcal{W}_{\text{social}} = \{\mathcal{W}_{\text{self}}, \mathcal{W}_{\text{others}}, \mathcal{W}_{\text{interaction}}\}
\end{equation}

\textbf{Other Agents' Models}:
\begin{equation}
\hat{a}_j^{(t+1)} = \pi_j(z_j^{(t)}, \mathcal{W}_{\text{others}})
\end{equation}

This enables:
\begin{itemize}
    \item Predicting other agents' actions
    \item Cooperative behavior
    \item Strategic reasoning
\end{itemize}
\end{frame}

\section{Advanced World Model Architectures}

\begin{frame}{Transformer-based World Models}
Recent work uses transformers for world modeling:

\begin{equation}
z_{t+1} = \text{Transformer}(z_{1:t}, a_{1:t})
\end{equation}

\textbf{Advantages}:
\begin{itemize}
    \item Long-range dependencies
    \item Parallel processing
    \item Attention mechanisms for relevant information
\end{itemize}

\textbf{Challenges}:
\begin{itemize}
    \item Computational complexity
    \item Memory requirements
    \item Training stability
\end{itemize}
\end{frame}

\begin{frame}{Diffusion World Models}
Inspired by diffusion models, we can formulate world models as:

\begin{equation}
p(z_{t+1} \mid z_t, a_t) = \int p(z_{t+1} \mid z_0, a_t) p(z_0 \mid z_t) dz_0
\end{equation}

where the diffusion process models the uncertainty in state transitions.

\textbf{Benefits}:
\begin{itemize}
    \item Natural uncertainty quantification
    \item Better handling of stochastic environments
    \item Improved sample efficiency
\end{itemize}
\end{frame}

\section{Challenges and Future Directions}

\begin{frame}{Current Challenges in World Models}
\begin{enumerate}
    \item \textbf{Model Complexity}: Balancing expressiveness and computational efficiency
    \item \textbf{Generalization}: Transferring knowledge across domains
    \item \textbf{Uncertainty}: Quantifying and propagating uncertainty
    \item \textbf{Scalability}: Scaling to complex, high-dimensional environments
\end{enumerate}

\textbf{Open Questions}:
\begin{itemize}
    \item How to learn world models from limited data?
    \item How to ensure world models are interpretable?
    \item How to combine multiple world models?
\end{itemize}
\end{frame}

\begin{frame}{Future Research Directions}
\begin{enumerate}
    \item \textbf{Neural-Symbolic World Models}: Combining neural networks with symbolic reasoning
    \item \textbf{Continual Learning}: Learning world models that adapt over time
    \item \textbf{Multi-modal World Models}: Integrating different sensory modalities
    \item \textbf{World Model Composition}: Combining multiple specialized world models
\end{enumerate}

\textbf{Long-term Vision}: World models that can:
\begin{itemize}
    \item Learn from minimal data
    \item Transfer knowledge across domains
    \item Reason about abstract concepts
    \item Collaborate with other agents
\end{itemize}
\end{frame}

\section{Applications and Case Studies}

\begin{frame}{World Models in Robotics}
\begin{itemize}
    \item \textbf{Sim-to-Real Transfer}: Train in simulation, deploy in real world
    \item \textbf{Manipulation Planning}: Plan complex manipulation tasks
    \item \textbf{Navigation}: Navigate in unknown environments
\end{itemize}

\textbf{Key Insight}: World models enable robots to reason about the physical world without direct interaction.
\end{frame}

\begin{frame}{World Models in Game AI}
\begin{itemize}
    \item \textbf{Strategy Games}: Model opponent behavior and game dynamics
    \item \textbf{Real-time Strategy}: Plan complex multi-agent strategies
    \item \textbf{Adversarial Games}: Model adversarial environments
\end{itemize}

\textbf{Examples}:
\begin{itemize}
    \item AlphaGo: Model-based planning in Go
    \item StarCraft AI: Strategic planning in RTS games
    \item Chess engines: Position evaluation and planning
\end{itemize}
\end{frame}

\begin{frame}{World Models in Scientific Discovery}
\begin{itemize}
    \item \textbf{Physics Simulation}: Model physical laws and phenomena
    \item \textbf{Chemical Reactions}: Predict reaction outcomes
    \item \textbf{Biological Systems}: Model complex biological processes
\end{itemize}

\textbf{Potential Impact}: World models could accelerate scientific discovery by enabling:
\begin{itemize}
    \item Virtual experimentation
    \item Hypothesis generation
    \item Theory validation
\end{itemize}
\end{frame}

\section{Theoretical Foundations}

\begin{frame}{Information-Theoretic Perspective}
World models can be understood through information theory:

\begin{equation}
I(S; Z) = H(S) - H(S \mid Z)
\end{equation}

where $I(S; Z)$ is the mutual information between observations $S$ and latent states $Z$.

\textbf{Optimal World Model}: Maximizes mutual information while minimizing complexity:
\begin{equation}
\max_{\theta} I(S; Z_\theta) - \lambda \text{Complexity}(Z_\theta)
\end{equation}
\end{frame}

\begin{frame}{Causal World Models}
For AGI, world models must capture causal relationships:

\begin{equation}
p(\text{effect} \mid \text{cause}) \neq p(\text{effect} \mid \text{correlation})
\end{equation}

\textbf{Causal Discovery}: Learning causal graphs from data
\begin{equation}
\mathcal{G} = \{(V, E) : V \text{ variables}, E \text{ causal edges}\}
\end{equation}

\textbf{Intervention}: Predicting effects of interventions
\begin{equation}
p(Y \mid \text{do}(X = x)) \neq p(Y \mid X = x)
\end{equation}
\end{frame}

\begin{frame}{World Models and Consciousness}
Some researchers argue that world models are related to consciousness:

\begin{itemize}
    \item \textbf{Self-Model}: Model of the agent itself
    \item \textbf{Attention}: Focus on relevant aspects of the world
    \item \textbf{Memory}: Integration of past experiences
    \item \textbf{Prediction}: Anticipation of future events
\end{itemize}

\textbf{Open Question}: Can sophisticated world models lead to artificial consciousness?
\end{frame}

\section{Implementation Considerations}

\begin{frame}{Computational Requirements}
World models have significant computational requirements:

\begin{itemize}
    \item \textbf{Training}: Large-scale optimization problems
    \item \textbf{Inference}: Real-time prediction and planning
    \item \textbf{Memory}: Storing and retrieving world state
    \item \textbf{Communication}: Sharing world models between agents
\end{itemize}

\textbf{Optimization Strategies}:
\begin{itemize}
    \item Distributed training
    \item Model compression
    \item Efficient architectures
    \item Hardware acceleration
\end{itemize}
\end{frame}

\begin{frame}{Evaluation Metrics}
How do we evaluate world models?

\begin{enumerate}
    \item \textbf{Prediction Accuracy}: How well does the model predict future states?
    \item \textbf{Planning Performance}: How well can the model be used for planning?
    \item \textbf{Generalization}: How well does the model transfer to new domains?
    \item \textbf{Efficiency}: How computationally efficient is the model?
\end{enumerate}

\textbf{Benchmarks}:
\begin{itemize}
    \item Atari games
    \item MuJoCo environments
    \item Real-world robotics tasks
    \item Multi-agent environments
\end{itemize}
\end{frame}

\section{Conclusion and Future Work}

\begin{frame}{Summary}
World models are fundamental for achieving AGI:

\begin{itemize}
    \item \textbf{Mathematical Foundation}: Rigorous theoretical framework
    \item \textbf{Practical Applications}: Success in various domains
    \item \textbf{Research Challenges}: Many open problems remain
    \item \textbf{Future Potential}: Key to artificial general intelligence
\end{itemize}

\textbf{Key Takeaways}:
\begin{enumerate}
    \item World models enable efficient learning and planning
    \item They are crucial for generalization and transfer
    \item Current methods show promise but need improvement
    \item Future work should focus on scalability and robustness
\end{enumerate}
\end{frame}

\begin{frame}{Open Research Questions}
\begin{enumerate}[(1.)]
    \item How can we learn world models from limited data?
    \item How do we ensure world models are interpretable and trustworthy?
    \item How can we combine multiple world models effectively?
    \item How do we handle non-stationary environments?
    \item How can world models enable artificial consciousness?
\end{enumerate}

\textbf{Call for Research}: These questions represent exciting opportunities for future research in AI and AGI.
\end{frame}

\begin{frame}{Thank You}
\centering
\textbf{Questions and Discussion}

\vspace{2em}

World models represent a crucial step toward artificial general intelligence. The mathematical foundations are solid, but many challenges remain.

\vspace{1em}

\textbf{Contact}: \href{mailto:hpp@bimsa.cn}{hpp@bimsa.cn}

\vspace{1em}

\textbf{References}: 
\begin{itemize}
    \item Ha \& Schmidhuber (2018): World Models
    \item Hafner et al. (2019): Dreamer
    \item Chen et al. (2021): Decision Transformer
    \item And many more...
\end{itemize}
\end{frame}

\end{document}
